\uuid{EVk4}
\exo7id{7300}
\titre{exo7 7300}
\auteur{mourougane}
\organisation{exo7}
\datecreate{2021-08-10}
\isIndication{false}
\isCorrection{false}
\chapitre{Groupe, anneau, corps}
\sousChapitre{Groupe, sous-groupe}
\cpge{mpsi-mp2i.raisonnement-vocabulaire-ensembliste.b, pcsi.raisonnement-vocabulaire-ensembliste.b, ptsi.raisonnement-vocabulaire-ensembliste.b, mpsi-mp2i.structures-algebriques.a, mpsi-mp2i.structures-algebriques.c, bcpst1.logique-ensembles-raisonnement.b}
\module{Algèbre}
\niveau{L2}
\difficulte{}

\contenu{
\question{
Soit $E$ un ensemble, on note $\mathcal{P}(E)$ l'ensemble des parties de $E$. Si $A, B \in \mathcal{P}(E)$, on note $A\Delta B = (A \cup B) \setminus (A \cap B)$ la \emph{différence symétrique} de $A$ et $B$. Montrer que pour toute partie $A$ de $E$, $A^2 = A$ et que le triplet $(\mathcal{P}(E),\Delta,\cap)$ forme un anneau commutatif.
}
}
